% TYLCV review paper -- IEEEtran journal format.
% Compile with: pdflatex TYLCV_review_IEEE.tex (run twice). Works on Overleaf.
% References are listed in order of first citation, as IEEE requires.
% Entries marked "TODO" have incomplete metadata and must be completed before submission.
\documentclass[journal]{IEEEtran}
\usepackage{cite}
\usepackage{amsmath}
\usepackage{graphicx}
\usepackage{array}
\usepackage{url}
\usepackage{booktabs}
\usepackage{tabularx}

\begin{document}

\title{Tomato Yellow Leaf Curl Virus (TYLCV): Biology, Vector Interactions, Host Resistance and Emerging Management Strategies --- A Review}

\author{Ritesh~Dwivedi%
\thanks{R. Dwivedi is with the Department of Agricultural Sciences, Chandigarh, India (e-mail: riteshdwivedi033@gmail.com).}%
\thanks{Manuscript received November 2026.}}

\markboth{Review Paper, November 2026}{Dwivedi: TYLCV Review}

\maketitle

\begin{abstract}
Tomato yellow leaf curl virus (TYLCV; genus \textit{Begomovirus}, family \textit{Geminiviridae}) is among the most damaging pathogens of tomato (\textit{Solanum lycopersicum}) worldwide. It is transmitted in a persistent, circulative manner by the \textit{Bemisia tabaci} cryptic species complex and can cause severe yield loss in unprotected crops. This review summarizes recent work (2021--2026) on TYLCV genome function, vector transmission, global spread and recombination, host resistance (\textit{Ty} genes), vector management, and biotechnological interventions including CRISPR/Cas genome editing and RNA interference. Existing studies are compared by approach, strengths and limitations, and the main research gaps are identified: durability of resistance under high disease pressure and heat, incomplete understanding of the virus--vector interface, insecticide resistance in whiteflies, and limited field validation of genome-edited and RNA-based protection. Durable control will require pyramiding of \textit{Ty} loci, vector exclusion and field-tested editing strategies within one integrated management framework.
\end{abstract}

\begin{IEEEkeywords}
TYLCV, begomovirus, \textit{Bemisia tabaci}, Ty genes, resistance breakdown, CRISPR/Cas, RNA interference, integrated pest management.
\end{IEEEkeywords}

\section{Introduction}
\IEEEPARstart{T}{omato} is one of the most widely grown vegetable crops, and its production is constrained by many pathogens. Among them, tomato yellow leaf curl disease (TYLCD) ranks among the most important, with the monopartite begomovirus TYLCV as its best-known causal agent \cite{yan2021}. Typical symptoms are upward leaf curling, marginal chlorosis, severe stunting and flower drop; infection before flowering can cause near-total loss of marketable fruit.

TYLCV was first associated with disease in the Middle East, and phylogeographic analyses place the origin of the first TYLCVs in that region, with later worldwide dissemination driven largely by trade and by invasive whitefly lineages \cite{marchant2023}. The disease now occurs across tropical, subtropical and warm-temperate production areas. This review (i) outlines TYLCV biology, (ii) synthesizes recent literature on transmission, evolution and management, (iii) compares the main control strategies, and (iv) identifies research gaps and future directions.

\section{Background and Related Work}
\subsection{Taxonomy and Genome Organization}
TYLCV belongs to the genus \textit{Begomovirus}; the family \textit{Geminiviridae} has recently been reorganized with the establishment of new genera \cite{roumagnac2022}, and begomoviruses are widely discussed as a highly successful and adaptable virus group \cite{fiallo2023}. Virions are twinned (geminate) particles formed by two incomplete icosahedra of about $20\times30$~nm enclosing a circular single-stranded DNA genome of roughly 2.8~kb. The genome carries six overlapping open reading frames (ORFs) on the virion-sense and complementary-sense strands, separated by an intergenic region containing the origin of replication (Table~\ref{tab:orf}).

\begin{table}[!t]
\caption{TYLCV Open Reading Frames and Functions}
\label{tab:orf}
\centering
\small
\begin{tabularx}{\columnwidth}{@{}l l X@{}}
\toprule
ORF & Protein & Main function \\
\midrule
V1 & Coat protein (CP) & Encapsidation, systemic movement, vector interaction \\
V2 & Pre-coat protein & Movement; strong silencing suppressor \\
C1 & Rep & Rolling-circle replication initiation \\
C2 & TrAP & Transcription control; silencing suppression \\
C3 & REn & Enhances replication \\
C4 & Pathogenicity factor & Symptom induction, movement, silencing suppression \\
\bottomrule
\end{tabularx}
\end{table}

Comparative work on silencing suppressors of TYLCD-associated viruses has shown that V2 is typically the strongest suppressor, with strain- and host-dependent activity of C2 and C4; this underlies much of the virus's ability to overcome RNA-silencing defenses.

\subsection{Host Range}
Besides tomato, TYLCV infects other solanaceous crops (pepper, eggplant, tobacco), common bean, and weeds such as \textit{Datura stramonium}, \textit{Solanum nigrum} and \textit{Malva parviflora}, which act as inoculum reservoirs.

\subsection{Vector and Transmission}
TYLCV is transmitted by \textit{B. tabaci} (notably MEAM1 and MED) in a persistent, circulative manner: virions are acquired from phloem, cross the midgut into the hemolymph, and reach the salivary glands before inoculation \cite{circ2021}. The whitefly endosymbiont \textit{Hamiltonella} GroEL chaperonin has been linked to protection of virions in the hemolymph and to differences in transmission efficiency between whitefly species. Whitefly salivary effectors also modulate plant hormone signaling, generally favoring salicylic-acid over jasmonic-acid defenses, which can benefit both the insect and the virus \cite{effectors2021}.

\section{Literature Review}
\subsection{Global Status and Epidemiology}
Recent global surveys emphasize that TYLCD now involves multiple virus species and strains and threatens every producing continent \cite{yan2021}. A genome-scale analysis of more than 800 TYLCV genomes reported that the Israel (IL) and Mild (Mld) strains occur across Africa, the Americas, Asia, Oceania and Europe, whereas recombinant strains were found mainly in the Middle East \cite{marchant2023}. In the western Mediterranean, a survey of field samples collected between 2015 and 2019 found the invasive recombinant TYLCV-IS76 in Morocco but not in Italy, showing country-specific shifts in TYLCV/TYLCSaV recombinant populations \cite{med2023}.

\subsection{Host Resistance and the Ty Loci}
Six major loci (\textit{Ty-1} to \textit{Ty-6}), introgressed from wild tomato relatives, have been used in breeding \cite{yan2021,elsappah2022}:
\begin{itemize}
\item \textit{Ty-1/Ty-3} are allelic and encode an RNA-dependent RNA polymerase that enhances methylation of viral DNA, giving broad geminivirus resistance \cite{elsappah2022}.
\item \textit{Ty-2} (chromosome 11, \textit{S. habrochaites}) encodes an NB-LRR protein and is effective against TYLCV but not all other begomoviruses \cite{elsappah2022}.
\item \textit{ty-5} is recessive and results from loss of function of the mRNA-surveillance factor Pelota (\textit{SlPelo}).
\item \textit{Ty-6} (chromosome 10) was recently identified as a variant of DNA polymerase delta subunit 1 (SlPOLD1) \cite{ty6_2024}; combined with \textit{Ty-3} or \textit{ty-5} it gives the highest resistance levels.
\item \textit{Ty-4} is a minor-effect locus.
\end{itemize}
Resistance is not absolute. Virus accumulation can still be detected in resistant plants, and breakdown of \textit{Ty-1}-mediated resistance has been reported in Spain, Morocco and Italy, linked to TYLCV/TYLCSV recombinants and high inoculum pressure \cite{elsappah2022,temp2021}. \textit{Ty-1/Ty-3} activity has also been reported to decline at higher temperature, whereas \textit{Ty-2} inhibition is more stable \cite{temp2021}.

\subsection{Vector Management}
Whitefly control combines exclusion (fine-mesh nets), reflective mulches, trapping and insecticides. Insecticide resistance is a growing problem: surveys of \textit{B. tabaci} on vegetable crops report resistance mutations across several chemical classes and cite hundreds of resistance cases globally \cite{insects2024}. Recent field trials show that rotations of synthetic insecticides still reduce whitefly numbers more consistently than biopesticides \cite{florida2023}, arguing for integrated and resistance-management programs. Biological control with predators, parasitoids and entomopathogens is increasingly reviewed as an IPM component \cite{biocontrol2025}, and combining resistant cultivars with exclusion has been associated with lower insecticide use and better biocontrol.

\subsection{Biotechnological Approaches}
\textit{CRISPR/Cas.} Two strategies dominate: targeting the viral genome and editing host susceptibility genes. Editing \textit{SlPelo} in tomato restricted viral DNA accumulation and conferred TYLCV resistance \cite{pramanik2021}. Targeting viral coding regions can generate escape variants, and CRISPR-mediated resistance has been extended to related begomoviruses such as ToLCNDV \cite{tolcndv2025}.

\textit{RNA interference.} Topically applied dsRNA from TYLCV C4 and V2 moved systemically in tomato and persisted for weeks \cite{tabein2022}, yet a controlled comparison reported no significant reduction of TYLCV-IL accumulation by external dsRNA, whereas a hairpin (hpRep) construct conferred strong resistance \cite{dsrna2021}. Nanoformulated dsRNA targeting both whitefly and virus is under development.

\section{Comparative Analysis of Existing Research}
Table~\ref{tab:comp} compares the main strategies.

\begin{table*}[!t]
\caption{Comparison of TYLCV Management Strategies}
\label{tab:comp}
\centering
\small
\begin{tabularx}{\textwidth}{@{}l l X X@{}}
\toprule
Strategy & Evidence & Strengths & Limitations \\
\midrule
\textit{Ty-1/Ty-3} breeding & \cite{elsappah2022,temp2021} & Broad geminivirus resistance; widely deployed & Breakdown under high pressure and recombinants; weaker at high temperature \\
\textit{Ty-2} & \cite{elsappah2022,temp2021} & Stable at higher temperature & Narrower virus spectrum \\
\textit{ty-5} / \textit{Ty-6} & \cite{pramanik2021,ty6_2024} & New mechanisms; best when stacked & Few multi-environment durability data \\
CRISPR (host gene) & \cite{pramanik2021} & Non-transgenic potential; restricts viral DNA & Regulation; pleiotropy of susceptibility-gene edits \\
CRISPR (viral target) & \cite{pramanik2021,tolcndv2025} & Heritable interference & Escape variants; delivery \\
Topical dsRNA/RNAi & \cite{tabein2022,dsrna2021} & Non-transgenic, rapid & Inconsistent protection; stability; cost \\
Insecticides & \cite{insects2024,florida2023} & Rapid vector suppression & Resistance; environmental load \\
Biological control & \cite{biocontrol2025} & Sustainable; IPM-compatible & Variable efficacy; transmission after brief feeding \\
Exclusion/mulches & \cite{circ2021} & No chemical input & Cost; ventilation; short protection window \\
\bottomrule
\end{tabularx}
\end{table*}

\section{Research Gaps}
\begin{enumerate}
\item \textit{Durability of resistance.} Few studies evaluate stacked \textit{Ty} loci across seasons, temperatures and virus populations, although recombinants are known to overcome \textit{Ty-1} \cite{elsappah2022,temp2021,med2023}.
\item \textit{Virus--vector interface.} Candidate whitefly proteins and effectors are described, but their specific roles in TYLCV transmission remain unresolved \cite{circ2021,effectors2021}.
\item \textit{Climate effects.} Temperature affects \textit{Ty} performance and likely vector dynamics; modeling suggests warmer conditions worsen outbreak persistence \cite{model2026}, but multi-year field data are scarce.
\item \textit{Field validation} of editing and RNA tools is limited \cite{pramanik2021,tabein2022,dsrna2021}.
\item \textit{Surveillance} of recombinants is geographically uneven \cite{marchant2023,med2023}.
\end{enumerate}

\section{Challenges and Limitations}
Resistance is partial and viral replication is not fully blocked, allowing selection of adapted variants. Insecticide resistance erodes chemical options \cite{insects2024,florida2023}. Regulatory and acceptance barriers affect genome-edited cultivars, and net-house costs limit use by smallholders. Research is concentrated on a few strains (IL, Mld), leaving other regions and recombinants poorly characterized.

\section{Future Research Directions}
Priorities include: (i) pyramiding \textit{Ty-1/Ty-3}, \textit{Ty-2}, \textit{ty-5} and \textit{Ty-6} through marker-assisted selection, with multi-location and multi-temperature testing \cite{elsappah2022,ty6_2024}; (ii) editing \textit{SlPelo} and multiplexed viral targets with field trials and off-target assessment \cite{pramanik2021}; (iii) dual-target dsRNA/nanocarrier systems against virus and vector; (iv) functional studies of whitefly effectors and endosymbiont proteins as targets for blocking transmission \cite{circ2021,effectors2021}; (v) genomic surveillance for recombinants \cite{marchant2023,med2023}; and (vi) integration of resistant cultivars, exclusion and biocontrol with insecticide-resistance management \cite{biocontrol2025,florida2023}.

\section{Conclusion}
TYLCV remains a major constraint on tomato production because of efficient whitefly transmission, rapid evolution and recombination, and the incomplete durability of single resistance sources. New resistance mechanisms (Pelota, POLD1), genome-editing demonstrations and improved IPM knowledge provide a stronger toolkit, but field-validated, multi-layered strategies are still needed. Stacked host resistance, vector exclusion and biological control, supported by genomic surveillance, offer the most realistic route to durable protection.

\begin{thebibliography}{99}
% TODO: Entries marked TODO lack author lists / exact titles; complete from the publisher page before submission.

\bibitem{yan2021} Z. Yan \textit{et al.}, ``The global dimension of tomato yellow leaf curl disease: Current status and breeding perspectives,'' \textit{Microorganisms}, vol. 9, no. 4, p. 740, 2021. [Online]. Available: \url{https://www.ncbi.nlm.nih.gov/pmc/articles/PMC8066563/} % TODO verify authors

\bibitem{marchant2023} W. G. Marchant \textit{et al.}, ``Global genomic analysis of tomato yellow leaf curl virus,'' \textit{Front. Virol.}, vol. 3, Art. no. 1221156, 2023. [Online]. Available: \url{https://www.frontiersin.org/articles/10.3389/fviro.2023.1221156/full} % TODO exact title, authors

\bibitem{roumagnac2022} P. Roumagnac \textit{et al.}, ``Establishment of five new genera in the family \textit{Geminiviridae},'' \textit{Arch. Virol.}, vol. 167, pp. 695--710, 2022. % TODO verify (from memory)

\bibitem{fiallo2023} E. Fiallo-Olive and J. Navas-Castillo, ``Begomoviruses: What is the secret(s) of their success?'' \textit{Trends Plant Sci.}, vol. 28, pp. 715--727, 2023. % TODO verify (from memory)

\bibitem{circ2021} [TODO: authors], ``[TODO: title; review of circulative begomovirus transmission by \textit{Bemisia tabaci}],'' \textit{Viruses}, vol. 13, no. 9, Art. no. 1808, 2021.

\bibitem{effectors2021} [TODO: authors], ``[TODO: title; review of whitefly salivary effectors],'' \textit{Front. Plant Sci.}, vol. 12, Art. no. 661141, 2021. [Online]. Available: \url{https://www.frontiersin.org/journals/plant-science/articles/10.3389/fpls.2021.661141/}

\bibitem{med2023} [TODO: authors], ``[TODO: title; western Mediterranean TYLCV/TYLCSaV recombinant survey, Morocco/Italy],'' \textit{Peer Community J.}, [TODO: year, vol.]. [Online]. Available: \url{https://peercommunityjournal.org/articles/10.24072/pcjournal.628/}

\bibitem{elsappah2022} A. H. El-Sappah \textit{et al.}, ``Natural resistance of tomato plants to tomato yellow leaf curl virus,'' \textit{Front. Plant Sci.}, vol. 13, Art. no. 1081549, 2022. [Online]. Available: \url{https://www.frontiersin.org/journals/plant-science/articles/10.3389/fpls.2022.1081549/}

\bibitem{ty6_2024} [TODO: authors], ``[TODO: title; identification of \textit{Ty-6} as SlPOLD1],'' \textit{Theor. Appl. Genet.}, 2024, doi: 10.1007/s00122-024-04803-w.

\bibitem{temp2021} [TODO: authors], ``[TODO: title; tomato begomovirus resistance, temperature and breakdown],'' \textit{Front. Plant Sci.}, vol. 12, Art. no. 685382, 2021. [Online]. Available: \url{https://www.frontiersin.org/journals/plant-science/articles/10.3389/fpls.2021.685382/}

\bibitem{insects2024} [TODO: authors], ``Insecticide resistance status and vector potential of \textit{Bemisia tabaci} populations on vegetable crops in T\"urkiye,'' \textit{Insects}, vol. 15, no. 6, Art. no. 438, 2024.

\bibitem{florida2023} [TODO: authors], ``[TODO: title; synthetic vs. biopesticide insecticide rotations for whitefly in Florida tomato, 2023 trials],'' \textit{Preprints.org}, 2024. [Online]. Available: \url{https://www.preprints.org/manuscript/202405.1136/v1}

\bibitem{biocontrol2025} [TODO: authors], ``[TODO: title; review of biological control of whiteflies],'' \textit{Front. Insect Sci.}, 2025. [Online]. Available: \url{https://www.frontiersin.org/journals/insect-science/articles/10.3389/finsc.2025.1684672/full}

\bibitem{pramanik2021} D. Pramanik \textit{et al.}, ``CRISPR/Cas9-mediated generation of pathogen-resistant tomato against tomato yellow leaf curl virus and powdery mildew,'' \textit{Int. J. Mol. Sci.}, 2021. [Online]. Available: \url{https://www.ncbi.nlm.nih.gov/pmc/articles/PMC7917697/} % TODO vol/issue/pages

\bibitem{tolcndv2025} [TODO: authors], ``[TODO: title; CRISPR/Cas9 targeting of tomato leaf curl New Delhi virus],'' \textit{J. Phytopathol.}, 2025.

\bibitem{tabein2022} [TODO: authors], ``[TODO: title; topical dsRNA from C4 and V2 against TYLCV-Mld in tomato],'' \textit{Biol. Plant.}, 2021. [Online]. Available: \url{https://www.bp.ueb.cas.cz/pdfs/bpl/2021/01/28.pdf}

\bibitem{dsrna2021} S. Tabein \textit{et al.}, ``[TODO: title; non-transgenic dsRNA vs. hairpin-RNA approaches against TYLCV-IL],'' \textit{J. Plant Prot.} (Shahid Chamran Univ. Ahvaz), 2022. [Online]. Available: \url{https://plantprotection.scu.ac.ir/article_17177.html}

\bibitem{model2026} [TODO: authors], ``[TODO: title; mathematical model of temperature-dependent TYLCV outbreak dynamics],'' \textit{J. Appl. Math.}, [TODO: year confirm]. [Online]. Available: \url{https://ideas.repec.org/a/hin/jnljam/3408574.html}

\end{thebibliography}

\end{document}
